
d'une
démonstration de la formule analogue dans le cas analytique complexe (ou rigide-analytique !). Rappelons que dans le cas analytique complexe, seul le cas d'un faisceau localement libre $F$ sur une variété non singulière compacte $X$, au-dessus d'un $Y$ ponctuel, est actuellement connu (grâce à la formule de l'index de Atiyah--Singer), le cas où $F$ est seulement supposé cohérent semblant échapper encore aux techniques connues jusqu'à présent.

\textbf{2.2.} Comme ingrédient essentiel de la formule de Riemann--Roch, rappelons aussi pour mémoire le \uline{problème de finitude} soulevé dans Exp. III: si $f:X\to Y$ est un morphisme propre, et $L_\bullet$ un complexe sur $X$ qui est pseudo-cohérent relativement à $Y$, alors $Rf_*(L_\bullet)$ est-il pseudo-cohérent? Une démonstration générale de cette conjecture mettra probablement en oeuvre des idées algébriques qui permettront de donner une démonstration du ``théorème'' de finitude de Grauert [7] sous des conditions notablement plus générales que dans loc. cit., en se plaçant d'emblée sur des espaces (ou topos) de base topologiquement annelés en $\C$-algébres $Y$ plus généraux que des espéces analytiques (englobant le cas où $Y$ serait un espace localement compact avec le faisceau des fonctions continues complexes, ou une variété $C^m$ avec le faisceau des fonctions complexes $C^m$, ou un espace analytique réel, etc).

\textbf{2.3.} Mentionnons enfin qu'il y aurait lieu de remplacer dans ces questions l'hypothèse de propreté sur le morphisme $f$ par l'hypothèse que les faisceaux de cohomologie $H^i(L_\bullet)$ des complexes envisagés sur $X$ aient leurs groupes propres sur $Y$.

\section*{3. Formule de Riemann--Roch ``sans dénominateurs'' pour une immersion}

\textbf{3.1.} Soit $f:Y\to X$ une immersion fermée régulière (Exp. VII), de faisceau conormal $\underline N$. Supposons pour simplifier que $X$ admet un faisceau inversible ample. La formule de Riemann--Roch pour $y\in K^\bullet(Y)$~:
\[
 \ch_X(f_*(y))=f_*\bigl(\ch_Y(y)\operatorname{Todd}(-\check{\underline N})\bigr),
\]
jointe à la formule générale (Exp. VII)~:
\[
 f_*(y')\,f_*(y'')=f\bigl(y'y''c_d(\underline N)\bigr),\qquad y',y''\in \operatorname{Gr}(Y)_{\underline Q},
\]
(en supposant $f$ de codimension $d$), permet d'exprimer les classes de Chern de $f_*(y)$ dans $\operatorname{Gr}^\bullet(X)_\Q$ sous la forme
